\documentclass[10pt]{../datasheet}

\usepackage[utf8]{inputenc}
\usepackage[english]{babel}
\usepackage[english]{isodate}
\usepackage{graphicx}
\usepackage[numbers,square]{natbib}
\usepackage{anysize}
\usepackage{graphicx} 
\marginsize{1.5cm}{1.5cm}{1cm}{1cm}
\usepackage{fancyhdr}
\renewcommand{\headrulewidth}{0pt}
\usepackage[dvipsnames]{xcolor}
\usepackage{hyperref}
\hypersetup{
    colorlinks=true,
    linkcolor=black,
    citecolor=CadetBlue,
    filecolor=CadetBlue,      
    urlcolor=CadetBlue,
}
\usepackage{multicol}
\setlength\columnsep{18pt}
\renewenvironment{abstract}
 {\par\noindent\textbf{\abstractname}\ \ignorespaces \\}
 {\par\noindent\medskip}
\renewcommand*{\thefootnote}{\fnsymbol{footnote}}
\usepackage{amsmath}


\usepackage{xcolor}
\usepackage{comment}
\usepackage{multirow}
\usepackage{float}
\usepackage{gensymb}
\usepackage{braket}
\usepackage{longtable}
\usepackage{caption}
\usepackage{longtable}
\usepackage{physics}
\usepackage{amsthm}
\usepackage{amsmath}
\usepackage{amssymb}
\usepackage{xcolor}
\usepackage{soul}
\usepackage{ marvosym }
\usepackage{wasysym}
\usepackage [autostyle, english = american]{csquotes}
\usepackage{fancyheadings}
\usepackage{enumitem}
\usepackage{multicol}


\MakeOuterQuote{"}

\newcommand{\xhat}{\hat{\vb x}}
\newcommand{\yhat}{\hat{\vb y}}
\newcommand{\zhat}{\hat{\vb z}}
\newcommand{\nhat}{\hat{\vb n}}
\newcommand{\rhat}{\hat{\vb r}}
\newcommand{\thetahat}{\hat{\vb \theta}}
\newcommand{\phihat}{\hat{\vb \phi}}


\renewcommand{\a}{\alpha}
\renewcommand{\b}{\beta}
\newcommand{\g}{\gamma}
\renewcommand{\d}{\delta}
\newcommand{\e}{\varepsilon}
\renewcommand{\l}{\lambda}
\renewcommand{\t}{\tau}
\newcommand{\w}{\omega}
\newcommand{\A}{\Alpha}
\newcommand{\B}{\Beta}
\newcommand{\G}{\Gamma}
\newcommand{\D}{\Delta}
\newcommand{\W}{\Omega}

\newcommand{\R}{\mathbb{R}}





\let\cite\citep
\usepackage{csvsimple-l3}
\usepackage{pdflscape}

\title{Electrical BUS Subsystem Literature Review}
\author{UTAT Space Systems}
\date{July 2026}
\revision{Revision 1}
\companylogo{\includegraphics[scale=0.022]{../figures/utat_ss.png}}

\begin{document}
\maketitle
\section{Overview of the Electrical BUS Subsystem}

The Electrical BUS subsystem's goal is to provide the CubeSat with a centralized
point of control, which at a high level will involve observing indicators and
driving actuators. More specifically, Electrical BUS will interact with the
PAYload, RF, and EPS subsystems to manage the state of affairs on the CubeSat.
Operations of the electrical subsystem may include directing the PAYload
subsystem to send data to the RF subsystem for downlinking, and acquiring uplink
data from RF to perform a reboot.


\section{Literature Review Methodology}

As it pertains to the design, the main components that make up the electrical
subsystem essentially boil down to three categories: control unit (that is,
device(s) responsible for carrying out generic operations of the electrical
subsystem, for example an MCU), protocol (that is, an interface for connecting
two or more devices and allowing them to share data in one way or another, for
example I2C), and peripherals (that is, device(s) responsible for a particular
function within the Electrical BUS subsystem, for example a watchdog). This
classification is not a recognized one, but has been found to be useful when
conducting this literature review since the overwhelming majority of the designs
can be deconstructed into components each of which would fit into one of these
categories [subjective].

In terms of sources, two major categories are surveyed:
\begin{itemize}
  \item Commercial OBCs discussed in NASA's Small-Sat State of the Art reports from 2022 to 2026 (\cite[pp.~215--226]{nasa_soa_2022}, \cite[pp.~225--237]{nasa_soa_2023}, \cite[pp.~227--239]{nasa_soa_2024}, \cite[pp.~241--253]{nasa_soa_2026}), with trends also observed.
  \item OBCs built by design teams around Canada and beyond.
\end{itemize}

\onecolumn

\section{Design Space}

This section outlines the resulting design space for the Electrical BUS
Subsystem based on the corresponding literature review. We establish a space of
common control units, protocols and periptherals for our internal design
process.

\section{Control Units}

The control-unit design space ranges from low-power microcontrollers to
heterogeneous processor--FPGA systems. NASA's State of the Art reports show
that most conventional CubeSat OBCs use microcontrollers or embedded
microprocessors because command execution, telemetry collection, timekeeping,
fault management, and subsystem coordination do not require high computational
performance \cite{nasa_soa_2022,nasa_soa_2023,nasa_soa_2024,nasa_soa_2026}.
Higher-performance FPGAs, application processors, and system-on-chip devices
are generally introduced when the OBC must also support image processing,
autonomy, high-rate payload data, or machine-learning workloads.

\begin{longtable}{|p{2.8cm}|p{4.6cm}|p{5.3cm}|p{3.8cm}|}
\caption{Candidate OBC control-unit architectures} \label{tab:control_architectures} \\ \hline
\textbf{Architecture} &
\textbf{Characteristics} &
\textbf{NASA SOA Examples} &
\textbf{Academic Examples} \\ \hline
\endfirsthead
\caption[]{Candidate OBC control-unit architectures (Continued)} \\ \hline
\textbf{Architecture} &
\textbf{Characteristics} &
\textbf{NASA SOA Examples} &
\textbf{Academic Examples} \\ \hline
\endhead
\hline
\endfoot

Single MCU or processor &
One device performs command execution, telemetry, timekeeping, fault
management, and subsystem control. Lowest complexity and generally lowest
power. &
EnduroSat OBC, ARM Cortex-M7
\cite[p.~229]{nasa_soa_2023};
AAC Clyde Space Kryten-M3, Cortex-M3
\cite[p.~228]{nasa_soa_2023}. &
STM32L4 PC/104 OBC; Athena TMS570; ORCASat TMS570
\cite{stm32l4_vst104_pub,athena_obc_site,orcasat_obc_site}. \\ \hline

Multicore SoC &
Multiple cores share one device and memory system. Provides additional
performance or lockstep capability, but does not remove the SoC as a
single point of failure. &
Argotec FERMI, dual-core LEON3FT
\cite[p.~229]{nasa_soa_2023};
SPiN MA61C, dual-core GR712RC
\cite[p.~236]{nasa_soa_2024};
Xiphos Q7S, dual-core Cortex-A9
\cite[p.~238]{nasa_soa_2024}. &
Athena and ORCASat use lockstep-capable ARM safety microcontrollers
\cite{athena_obc_ti_prod,orcasat_obc_ds}. \\ \hline

Dual redundant processors &
Two processor devices provide cold, warm, or active redundancy. One processor
may take control following a fault in the primary device. &
Andes OBC, two Cortex-M7 processors with full cold redundancy
\cite[p.~230]{nasa_soa_2026};
Bradford Binary OBC, dual Cortex-R5 with internal cold redundancy
\cite[p.~229]{nasa_soa_2023};
Bradford Stellar OBC, dual PolarFire SoCs with internal cold redundancy
\cite[p.~230]{nasa_soa_2026}. &
No directly equivalent two-device redundant architecture was identified in
the academic survey. \\ \hline

Controller plus FPGA or DPU &
A low-power controller performs critical spacecraft control while an FPGA or
higher-performance processor handles payload data or acceleration. The
secondary unit may be powered down when unused. &
KP Labs Antelope, RM57 OBC plus Zynq UltraScale+ DPU
\cite[p.~230]{nasa_soa_2023};
NOVI Gen~2, Vorago MCU plus Versal AI Edge SoC
\cite[p.~232]{nasa_soa_2026}. &
NEUDOSE, primary AVR32 plus secondary Xilinx Zynq-7000
\cite{neudose_obc_pub}. \\ \hline

FPGA-centric &
Control or processing is implemented using programmable logic and,
optionally, a soft processor. Suitable for deterministic parallel interfaces
and high-rate data paths, but more difficult to develop as a general-purpose
OBC. &
LASP Generic SmallSat FPGA Board, Kintex-7
\cite[p.~230]{nasa_soa_2023};
KP Labs Lion, Kintex UltraScale
\cite[p.~230]{nasa_soa_2023}. &
No FPGA-only primary OBC was identified in the academic survey. \\ \hline

Distributed controllers &
Subsystem-local processors control the EPS, ADCS, communications, or payload
and exchange data with a coordinating OBC. This improves fault containment but
increases network and software complexity. &
NASA identifies federated, distributed, heterogeneous, and mixed-criticality
avionics as viable SmallSat architectures
\cite[p.~225]{nasa_soa_2023}
\cite[p.~221--222]{nasa_soa_2026}. &
Individual subsystem controllers used alongside the primary academic OBC. \\ \hline
\end{longtable}

All architectures will be considered during preliminary design; however, the
current baseline is a centralized microcontroller-based OBC with independent
watchdog and reset circuitry. This approach is the most common among the
academic systems surveyed and provides the lowest-risk balance of power,
determinism, software complexity, and flight heritage. An FPGA or secondary
processor may be added later if payload-processing requirements cannot be met
by the primary MCU.

\section{Protocols}

\section{Periptherals}

\section{COTS (Commercial Off the Shelf) Design Alternatives}

\textit{Description: Determine what COTS solutions exist if we cannot produce a certain part of the subsystem in house. Find a COTS option for each in house design.}

\textit{Will inform trades later on for COTS versus in house risks.}

This section outlines alternative off-the-shelf components that would in part or
in full fulfil the operating requirements of the subsystem.

\textbf{TODO}

\section{System Architecture and Budgets}

In this section we concern ourselves with the following considerations: power
draw, physical dimentions, radiation tolerance, and weight for the purposes of
the overall design and integration with other teams.

\subsection{Power draw}
Based on the commercial onboard-computing systems surveyed in NASA's
State of the Art reports
\cite{nasa_soa_2022,nasa_soa_2023,nasa_soa_2024,nasa_soa_2026},
OBC power consumption can be divided into three broad categories according
to platform complexity and computational capability. The reports do not use
a consistent average-power metric; reported values are variously identified
as idle, static, typical, active-mode, full-load, or peak consumption.
Consequently, Table~\ref{tab:power_trends} preserves the operating-state
terminology used by the manufacturers.

The 2022 edition identifies representative onboard-computing products but
does not tabulate their power consumption
\cite[Table~8-1, pp.~214--217]{nasa_soa_2022}. Quantitative power values are
therefore drawn primarily from the 2023, 2024, and 2026 editions
\cite[Table~8-1, pp.~228--233]{nasa_soa_2023}
\cite[Table~8-1, pp.~230--236]{nasa_soa_2024}
\cite[Table~8-4, pp.~229--234]{nasa_soa_2026}.

\begin{longtable}{|p{2.7cm}|p{3.2cm}|p{5.0cm}|p{5.4cm}|}
  \caption{OBC power-consumption trends and reported operating states} \label{tab:power_trends} \\ \hline
  \textbf{Platform Complexity}
   &
  \textbf{Reported Power Profile}
   &
  \textbf{Design Characteristics}
   &
  \textbf{Example Designs}
  \\ \hline
  \endfirsthead
  \caption[]{OBC power-consumption trends and reported operating states (Continued)} \\ \hline
  \textbf{Platform Complexity}
   &
  \textbf{Reported Power Profile}
   &
  \textbf{Design Characteristics}
   &
  \textbf{Example Designs}
  \\ \hline
  \endhead
  \hline
  \endfoot

  Low
  (Bare-Metal / Control)
   &
  \textbf{Minimum listed:} 0.05~W

  \textbf{Idle/steady:} approximately 0.2--0.5~W

  \textbf{Peak:} generally below 1~W
   &
  Simple command and data handling, housekeeping, telemetry, and fault
  management. These systems are normally based on low-power microcontrollers
  and execute bare-metal software or a small real-time operating system.
  Their power consumption is relatively insensitive to computational load.
   &
  Pumpkin PPM family:
  0.05~W, operating state not specified
  \cite[p.~231]{nasa_soa_2023}
  \cite[p.~235]{nasa_soa_2024}
  \cite[p.~233]{nasa_soa_2026}.

  AAC Clyde Space Kryten-M3:
  0.4~W
  \cite[p.~228]{nasa_soa_2023}
  \cite[p.~230]{nasa_soa_2024}
  \cite[Table~8-4]{nasa_soa_2026}.

  C3S Electronics OBC:
  0.46~W in test mode
  \cite[p.~229]{nasa_soa_2023}
  \cite[p.~232]{nasa_soa_2024}
  \cite[p.~230]{nasa_soa_2026}.

  SkyLabs NANOobc-2:
  $<0.9$~W idle and $<1$~W peak
  \cite[p.~232]{nasa_soa_2023}
  \cite[p.~236]{nasa_soa_2024}.
  \\ \hline

  Medium
  (Integrated OBC)
   &
  \textbf{Idle:} approximately 0.2--1~W

  \textbf{Nominal/steady:} approximately 1--3~W

  \textbf{Peak:} approximately 1.5--5~W
   &
  Typical integrated CubeSat OBCs combining command and data handling,
  mass-memory management, communication interfaces, fault protection, and
  basic payload management. Power may vary with processor utilization,
  interface activity, and memory access.
   &
  SPiN MA61C CubeSat:
  1--1.2~W
  \cite[p.~232]{nasa_soa_2023}
  \cite[p.~236]{nasa_soa_2024}
  \cite[Table~8-4]{nasa_soa_2026}.

  AAC Clyde Space Sirius OBC \& TCM:
  1.3~W
  \cite[p.~228]{nasa_soa_2023}
  \cite[p.~230]{nasa_soa_2024}
  \cite[Table~8-4]{nasa_soa_2026}.

  EnduroSat OBC I:
  0.2~W idle and 1.5~W peak
  \cite[p.~229]{nasa_soa_2023}
  \cite[Table~8-1]{nasa_soa_2024}
  \cite[p.~230]{nasa_soa_2026}.

  Resilient Computing RadPC:
  1.5~W
  \cite[p.~232]{nasa_soa_2023}
  \cite[p.~236]{nasa_soa_2024}
  \cite[p.~233]{nasa_soa_2026}.

  SatRev OBC:
  2.8~W typical and 5~W peak
  \cite[p.~233]{nasa_soa_2026}.
  \\ \hline

  High
  (Edge / AI / SoM)
   &
  \textbf{Idle/static:} approximately 5--15~W

  \textbf{Typical/active:} approximately 6--10~W

  \textbf{Peak/full load:} commonly 10--20~W, with larger payload
  computers reaching substantially higher values
   &
  Heterogeneous systems-on-module and payload-processing computers supporting
  image processing, software-defined radio, autonomous navigation, data
  compression, and machine-learning inference. Power is strongly dependent on
  processor utilization, FPGA configuration, accelerator activity, and
  payload duty cycle.
   &
  Novo Space GPU002AF:
  6--11~W in active mode
  \cite[p.~231]{nasa_soa_2023}
  \cite[p.~235]{nasa_soa_2024}
  \cite[p.~233]{nasa_soa_2026}.

  KP Labs Leopard DPU:
  7~W static and up to 10~W during deep-learning inference
  \cite[p.~230]{nasa_soa_2023}
  \cite[p.~234]{nasa_soa_2024}
  \cite[p.~231]{nasa_soa_2026}.

  Aitech S-A1760:
  $\leq5$~W idle, 8--10~W under a typical CUDA load, and
  $<20$~W at full utilization
  \cite[p.~228]{nasa_soa_2023}
  \cite[Table~8-1]{nasa_soa_2024}
  \cite[Table~8-4]{nasa_soa_2026}.

  EnduroSat GPC:
  $<15$~W idle and 130~W peak
  \cite[p.~231]{nasa_soa_2026}.
  \\ \hline
\end{longtable}

Based on these results, a nominal power-design envelope of 1--10~W is adopted
for the proposed OBC. This range encompasses conventional integrated CubeSat
OBCs and moderate edge-processing workloads, while transient loads above 10~W
are treated separately as peak operating conditions. The selected range is a
project design boundary rather than the full commercial-hardware envelope, which
extends from approximately 0.05~W for low-power microcontroller boards to 130~W
peak for high-performance payload computers.

[TODO: peak loads and durations]

\subsection{Dimentions}

\subsection{Weight}

\subsection{Radiation tolerance}

\clearpage
\begin{landscape}
  \section{Appendix A: Academic OBCs}

  \begin{longtable}{|p{2.6cm}|p{2.0cm}|p{3.8cm}|p{3.8cm}|p{1.4cm}|p{6.4cm}|p{1.0cm}|}
    \hline \textbf{Name}                                                                                                                                                                                                                                                              & \textbf{Manufacturer} & \textbf{MCU} & \textbf{Protocols} & \textbf{Standard} & \textbf{Notes} & \textbf{Ref} \\\hline \endhead
    \hline
    STM32L4 PC/104 academic OBC                                                                                                                                                                                                                                                       &
    VST104 thesis design                                                                                                                                                                                                                                                              &
    STM32L4—ARM-based                                                                                                                                                                                                                                                                 &
    CAN                                                                                                                                                                                                                                                                               &
    PC104                                                                                                                                                                                                                                                                             &
    pc 104 cubesat design.                                                                                                                                                                                                                                                            &
    \cite{stm32l4_vst104_pub}                                                                                                                                                                                                                                                                                                                                                                         \\ \hline
    Athena OBC                                                                                                                                                                                                                                                                        &
    University of Alberta (Albertasat)                                                                                                                                                                                                                                                &
    TI TMS570LC4357 (ARM Cortex R5F)                                                                                                                                                                                                                                                  &
                                                                                                                                                                                                                                                                                      &
                                                                                                                                                                                                                                                                                      &
    Embedded memory type: Flash. Features ideal for satellite OBC application, including radiation \& space reliability (rad-hard, SEL immune 48 MeV at 125\textdegree{}C, LAT 30krad), high computational performance (lockstep dual-core compare, 300MHz CPU clock, 1.66DMIPS/MHz). &
    \cite{athena_obc_site, athena_obc_ti_app, athena_obc_ti_prod}                                                                                                                                                                                                                                                                                                                                     \\ \hline
    Orcasat OBC                                                                                                                                                                                                                                                                       &
    University of Victoria                                                                                                                                                                                                                                                            &
    TI TMS570LS0714 (Dual ARM Cortex-R4F)                                                                                                                                                                                                                                             &
                                                                                                                                                                                                                                                                                      &
                                                                                                                                                                                                                                                                                      &
    Embedded memory type: Flash. Operates in lockstep, ECC integrated flash \& RAM banks, radiation-tested up to 6 krad, 80MHz core clock.                                                                                                                                            &
    \cite{orcasat_obc_site, orcasat_obc_ds}                                                                                                                                                                                                                                                                                                                                                           \\ \hline
    LORIS OBC                                                                                                                                                                                                                                                                         &
    Dalhousie University                                                                                                                                                                                                                                                              &
    NXP i.MX6DL (Single-core ARM Cortex-A9)                                                                                                                                                                                                                                           &
    USB 2.0; SPI; I2C; PWM; UART; GPIO; ADC                                                                                                                                                                                                                                           &
                                                                                                                                                                                                                                                                                      &
    Embedded memory type: Flash. Low profile, power efficient SoC SBC. 4GB eMMC flash storage, 256MB DDR3 RAM, 800MHz operational frequency, -40\textdegree{}C to 85\textdegree{}C temperature range.                                                                                 &
    \cite{loris_obc_pub, loris_obc_nxp}                                                                                                                                                                                                                                                                                                                                                               \\ \hline
    NEUDOSE OBC                                                                                                                                                                                                                                                                       &
    McMaster University                                                                                                                                                                                                                                                               &
    Primary: AVR32. Secondary: Xilinx Zynq-7000 SoC                                                                                                                                                                                                                                   &
                                                                                                                                                                                                                                                                                      &
                                                                                                                                                                                                                                                                                      &
    Primary is COTS chosen for high TRL to reduce mission risk; secondary is designed for radiation tolerance using rad-hard and non-rad-hard components.                                                                                                                             &
    \cite{neudose_obc_pub}                                                                                                                                                                                                                                                                                                                                                                            \\ \hline
    RADSAT-SK OBC                                                                                                                                                                                                                                                                     &
    University of Saskatchewan                                                                                                                                                                                                                                                        &
    ARM926EJ-S                                                                                                                                                                                                                                                                        &
                                                                                                                                                                                                                                                                                      &
                                                                                                                                                                                                                                                                                      &
    Embedded memory type: SDRAM. Uses ISISpace satellite-subsystems libraries, contains MMU, designed to interface with external SDRAM.                                                                                                                                               &
    \cite{radsat_sk_obc_git}                                                                                                                                                                                                                                                                                                                                                                          \\ \hline
  \end{longtable}
\end{landscape}
\clearpage

\begin{landscape}
  \section{Appendix B: Table of commercial OBCs}

  \begin{longtable}{|p{2.6cm}|p{2.0cm}|p{3.8cm}|p{3.8cm}|p{1.4cm}|p{6.4cm}|p{1.0cm}|}
    \hline \textbf{Name}                                                                                                             & \textbf{Manufacturer} & \textbf{MCU} & \textbf{Protocols} & \textbf{Standard} & \textbf{Notes} & \textbf{Ref} \\\hline \endhead
    \hline
    CFC-300                                                                                                                          &
    Innoflight                                                                                                                       &
    Xilinx Zynq SoC with dual-core ARM Cortex-A9 and FPGA fabric                                                                     &
    CAN; I2C; RS422; RS485; SPI; SpaceWire                                                                                           &
                                                                                                                                     &
    general-purpose; high performance. Fully on-orbit programmable; FPGA bit streams reloadable.                                     &
    \cite{cfc300_prod, cfc300_ds}, \cite[p.~222]{nasa_soa_2022}, \cite[p.~233]{nasa_soa_2023}, \cite[p.~235]{nasa_soa_2024}, \cite[p.~249]{nasa_soa_2026}                                                                                            \\ \hline
    CubeComputer V4.1                                                                                                                &
    CubeSpace                                                                                                                        &
    32-bit ARM Cortex-M3 based MCU                                                                                                   &
    CAN; I2C; SPI; UART                                                                                                              &
    PC104                                                                                                                            &
    General-purpose CubeSat OBC; Piggyback header—mission specific expansion. Used as ADCS OBC on QB50 precursor satellites.         &
    \cite{cubecomputer_v41_prod, cubecomputer_v41_ds}                                                                                                                                                                                                \\ \hline
    Gen 2 OBC                                                                                                                        &
    NOVI Space                                                                                                                       &
    Vorago ARM Cortex M4                                                                                                             &
                                                                                                                                     &
    PC104                                                                                                                            &
    Embedded memory type: FRAM.                                                                                                      &
    \cite[p.~222]{nasa_soa_2022}, \cite[p.~234]{nasa_soa_2023}, \cite[p.~236]{nasa_soa_2024}, \cite[p.~250]{nasa_soa_2026}                                                                                                                           \\ \hline
    iOBC                                                                                                                             &
    ISISPACE                                                                                                                         &
    Some Atmel ARM9 (32-bit ARM9) [datasheet]                                                                                        &
    GPIO; I2C; RS232; RS422; RS485; SPI; UART; USB                                                                                   &
                                                                                                                                     &
                                                                                                                                     &
    \cite{iobc_prod, iobc_ds}                                                                                                                                                                                                                        \\ \hline
    Main Bus Unit SatBus 3C2                                                                                                         &
    NanoAvionics                                                                                                                     &
                                                                                                                                     &
    CAN; I2C; SPI; UART; USB                                                                                                         &
                                                                                                                                     &
    No datasheet; includes ADCS.                                                                                                     &
    \cite{satbus_3c2_prod}                                                                                                                                                                                                                           \\ \hline
    NANOMIND A3200                                                                                                                   &
    GOMSPACE                                                                                                                         &
    Atmel AT32UC3C MCU / AVR32 MCU                                                                                                   &
    CAN; Cubesat Space Protocol; I2C; SPI; UART                                                                                      &
                                                                                                                                     &
    Tiny.                                                                                                                            &
    \cite{nanomind_a3200_prod, nanomind_a3200_ds}, \cite[p.~223]{nasa_soa_2022}, \cite[p.~234]{nasa_soa_2023}, \cite[p.~236]{nasa_soa_2024}, \cite[p.~250]{nasa_soa_2026}                                                                            \\ \hline
    NANOMIND HP MK3                                                                                                                  &
    GOMSPACE                                                                                                                         &
    Xilinx Zynq 7030 (FPGA + ARM) [product page]                                                                                     &
    CAN; RS422; RS485; SPI; SpaceWire; USB                                                                                           &
                                                                                                                                     &
    High performance (ARM dual core); FPGA included. Flew on SDR MK III; has PPS.                                                    &
    \cite{nanomind_hp_mk3_prod, nanomind_hp_mk3_ds}, \cite[p.~225]{nasa_soa_2022}, \cite[p.~236]{nasa_soa_2023}, \cite[p.~238]{nasa_soa_2024}, \cite[p.~252]{nasa_soa_2026}                                                                          \\ \hline
    NanoMind HP (PolarFire)                                                                                                          &
    GOMSPACE                                                                                                                         &
    32-bit RISC-V core / PolarFire                                                                                                   &
                                                                                                                                     &
    PC104                                                                                                                            &
    Embedded memory type: DDR3.                                                                                                      &
    \cite[p.~225]{nasa_soa_2022}, \cite[p.~236]{nasa_soa_2023}, \cite[p.~238]{nasa_soa_2024}, \cite[p.~252]{nasa_soa_2026}                                                                                                                           \\ \hline
    OBC - P4                                                                                                                         &
    Space Inventor                                                                                                                   &
    ARM Cortex-M7                                                                                                                    &
    CAN                                                                                                                              &
                                                                                                                                     &
    four OBC modules; hot/cold redundancy solutions (Datasheet: request on satsearch). FreeRTOS(tm) real time operating system.      &
    \cite{obc_p4_prod}                                                                                                                                                                                                                               \\ \hline
    Sirius OBC LEON3FT                                                                                                               &
    AAA Clyde Space                                                                                                                  &
    32-bit LEON3FT (IEEE-1754 SPARC v8) fault-tolerant processor                                                                     &
    CAN; GPIO; RS422; RS485; SpaceWire                                                                                               &
    SpaceWire 50 Mbps                                                                                                                &
    Reliable high-performance for advanced missions; designed and qualified for five years in LEO. Embedded memory type: NAND Flash. &
    \cite{sirius_obc_prod, sirius_obc_ds}, \cite[p.~219]{nasa_soa_2022}, \cite[p.~230]{nasa_soa_2023}, \cite[p.~232]{nasa_soa_2024}, \cite[p.~246]{nasa_soa_2026}                                                                                    \\ \hline
    [the] ONBOARD COMPUTER                                                                                                           &
    EnduroSat                                                                                                                        &
    ARM Cortex M7                                                                                                                    &
    I2C                                                                                                                              &
    PC104                                                                                                                            &
    Option to include GNSS receiver (Datasheet: [reg walled]). Embedded memory type: Flash.                                          &
    \cite{endurosat_obc_prod}, \cite[p.~222]{nasa_soa_2022}, \cite[p.~233]{nasa_soa_2023}, \cite[p.~235]{nasa_soa_2024}, \cite[p.~249]{nasa_soa_2026}                                                                                                \\ \hline
  \end{longtable}
\end{landscape}
\clearpage

\bibliographystyle{ieeetr}
\bibliography{references}

\end{document}
